% !TEX TS-program = pdflatexmk
\documentclass[12pt]{article}

% Layout.
\usepackage[top=1in, bottom=0.75in, left=1in, right=1in, headheight=1in, headsep=6pt]{geometry}

% Fonts.
\usepackage{mathptmx}
\usepackage[scaled=0.86]{helvet}
\renewcommand{\emph}[1]{\textsf{\textbf{#1}}}

% Misc packages.
\usepackage{amsmath,amssymb,latexsym}
\usepackage{graphicx,hyperref}
\usepackage{array}
\usepackage{xcolor}
\usepackage{multicol}
\usepackage{tabularx,colortbl}
\usepackage{enumitem}

\hypersetup{
    colorlinks=true,
    linkcolor=blue,
    filecolor=magenta,      
    urlcolor=blue,
    pdftitle={Calc I Syllabus},
    pdfpagemode=FullScreen,
    }

\def\mailto#1{\href{mailto:#1}{#1}}

% Paragraph spacing
\parindent 0pt
\parskip 6pt plus 1pt
\def\tableindent{\hskip 0.5 in}
\def\ts{\hskip 1.5 em}

\usepackage{fancyhdr}
\pagestyle{fancy} 
\lhead{\large\sf\textbf{MATH F251: Calculus I (asynchronous)}}
\rhead{\large\sf\textbf{Spring 2026 Syllabus}}

\newcommand{\localhead}[1]{\par\smallskip\textbf{#1}\nobreak\\}%
\def\heading#1{\localhead{\large\emph{#1}}}
\def\subheading#1{\localhead{\emph{#1}}}

\newenvironment{clist}%
{\bgroup\parskip 0pt\begin{list}{$\bullet$}{\partopsep 4pt\topsep 0pt\itemsep -2pt}}%
{\end{list}\egroup}%

\begin{document}

% \heading{Course Description}

\strut\par\vskip-12pt
\heading{Essential Information}

\vskip -12pt
\strut\hbox to \hsize{\tableindent\vtop{\halign{#\hfill\ts&#\hfil\cr
{\emph{Instructor}}&Dr. James Gossell (section 001) \cr
%\strut & \cr
{\emph{Email}}&jegossell@alaska.edu \cr
%\strut & \cr
{\emph{Physical Office}}&Chapman 301D \cr
\emph{Prerequisite} & MATH F151X and MATH F152X; or MATH F156X; or placement.\cr
\strut & \cr
{\emph{Required Text}} &\textit{OpenStax Calculus Volume 1} by G. Strang \& E. Herman,\cr
&\url{https://openstax.org/details/books/calculus-volume-1} \cr
&(optional print copy)\texttt{ISBN-13: 978-1938168024}\cr
  \strut & \cr
 {\emph{Required Technology}}&\textbullet \hspace{.05in} A scanner, smartphone, or camera with software or app for\cr
& \quad scanning documents and uploading them as PDFs\cr
& \textbullet \hspace{.05in} A printer or a tablet (e.g., iPad) where you can annotate documents\cr
& \textbullet \hspace{.05in} Reliable internet access \cr
\strut & \cr
  {\emph{Course Materials}}&Canvas (\url{https://www.uaf.edu/uaf/current/canvas.php})\cr
&Calculus I Webpage (\url{https://uaf-math.github.io/calc1/})\cr
  }
\hfil}}

\vskip -12pt
\heading{Description, Course Goals \& Student Learning Outcomes}
Calculus is one of mathematics' premiere computational tools.  It has
pervasive applications in all the sciences and is part of the UAF core 
curriculum.  The two principal tools of calculus are differentiation and integration. Differentiation concerns how changes in one variable affect another.
How does a population of bacteria change as time changes?  
How does the temperature of the ocean change as depth increases?
Integration, on the other hand, is a kind of reverse process to
differentiation. 

Students completing the course will have the mathematical foundation to be successful in Calculus II and other courses
requiring this background.  Specifically, students will
\begin{clist}
\item understand the role of limits in the definition of a derivative 
and be able to compute elementary derivatives from this definition,
\item understand the definition of a continuous function and identify
continuous/discontinuous functions,
\item develop the skills to compute standard derivatives,
\item be able to apply derivatives to common types of applied problems,
\item understand the definition of the definite integral,
\item be able to appply the Fundamental Theorem of Calculus to
compute definite integrals,
\item be able to apply integration to common types of applied problems.
\end{clist}

\heading{Time Commitment}
This is a 4-credit course, which means that a well-prepared student should expect to spend around 12
hours per week actually studying and doing work for this course. For students who may be missing substantial
prerequisite content knowledge, the time commitment is greater. The best way to manage a time
commitment of this magnitude is to schedule the hours into your day just as you would a paid job.

\heading{A Typical Week}
All the materials for this course are available online, linked either from Canvas or from the public Calculus I webpage. The course is organized as weekly modules in Canvas. Although this is an online course, \emph{it is not a self-paced course}. There are weekly
deadlines and a grading scheme that rewards timely completion. Within each module, tasks are
organized as daily chunks of work. The daily tasks are to be used as guides---you do not have
to do any task on any particular day with the exception of the two Midterms, the Derivative Proficiency and Derivative Proficiency Retake (optional), and the
Final Exam. However, there are regular deadlines, so you must complete tasks by specific days. Homework and quizzes may be completed in advance of the deadlines. 

Recommended typical workflow for this course is as follows: 
\begin{itemize}
\item {\bf Monday:}\\
Read the weekly announcements. Take a {\bf quiz} from the previous week's material and check your answers. After you have revised any errors in your quiz, upload it to Gradescope.
\item {\bf Tuesday:}\\
View the recitation video and corresponding worksheet to practice specific algebra or calculus skills needed for the assignments for that week. Complete any check-in assignments.
\item {\bf Wednesday-Friday:}\\
Read the book and/or watch videos for each section assigned that week. Work the homework problems for each section assigned that week and check your answers. After you have revised any errors in your homework, upload it to Gradescope.
\item {\bf Saturday-Sunday:}\\
Review the topics for the week and prepare to take the weekly quiz on Monday.
\end{itemize}

\vspace{.1in}

\heading{Tentative Schedule}
A day-to-day schedule is posted on Canvas and on the UAF Calculus I course website. 

This schedule mirrors the tasks for the in-person course. A student is free to organize their time as they choose. The hard deadlines are in
\textcolor{red}{red} in the row labeled Deadlines. These deadlines represent the \textbf{last} day to complete these tasks
and still be considered on time. The daily schedule is set up so that you can stay well ahead of
those deadlines. You should consult this schedule routinely. We may make minor adjustments to
the schedule, which will be announced in advance.

%We will use Canvas to send announcements. If we (your instructor/TA) need to contact you, we will send an email to you via Canvas. Thus, you will want to make sure that the email address in Canvas is one that you check regularly. Note that in Canvas it is possible to set up text alerts. However, you must login to Canvas and adjust the setting for your account. Neither email nor text alerts are automatic.

\vspace{.1in}

\heading{Online Course Materials}
All course materials can be accessed via Canvas. In addition, you will find a wealth of useful material at the public webpage: \href{https://uaf-math.github.io/calc1/}{https://uaf-math.github.io/calc1/}.

\vspace{.1in}

%\begin{table}[h]
%\centering
%\begin{tabular}{rl}
%where to find it&what you are looking for\\
%\hline \hline
%&syllabus and day-by-day schedule for the semester\\
%&ALEKS logistics for Weeks 1 and 2\\
%&homework problem sets\\
%public Calculus I webpage&practice quizzes, proficiencies, midterms, final exams\\
%\href{https://uaf-math251.github.io/}{https://uaf-math251.github.io/}&solutions to practice quizzes, proficiencies, midterms, final exam\\
%&videos\\
%&recitation worksheets with solutions\\
%&textbook\\
%\hline
%&announcements and reminders\\
%&link to Gradescope (to turn in homework)\\
%\href{https://www.uaf.edu/uaf/current/canvas.php}{Canvas}& complete solutions to homework\\
%& your grades\\
%&textbook\\
%\end{tabular}
%\caption{{\Large{\textbf{Online Course Materials}}}}
%\label{wheretofind}
%\end{table}



%\heading{Weeks 1 \& 2 Logistics}
%The first week of the course is devoted to prerequisite review. Consequently, the homework and quiz mechanics for the first two weeks are different from that of the remainder of the semester.
%
%Instead of the usual written homework, you will be
%working with a program called ALEKS to refresh precalculus skills. The first
%quiz will be an ALEKS-based assessment. (See \href{https://uaf-math251.github.io/week1-ALEKS.html}{Week 1 Details} for step-by-step instructions.)\\ 
%During the first weeks you will:
%\begin{clist}
%\item enroll in the UAF MATH 251 Spring 2022 Cohort of ALEKS\\
%class code: \textbf{VX3AW-T9LNW} %VX3AW-T9LNW
%\item complete an initial placement test (approx 1-2 hours) by \textbf{Tuesday August 30} at 11:59pm \\
%Finishing the initial placement on time is worth 10 points added to Quiz 1, your proctored ALEKS score.
%\item complete 90\% of the ALEKS pie \: OR \: spend 5 hours in Learning Mode by Monday, September 5 at 11:59pm which will count as your first homework grade.
%\item complete a \emph{proctored} ALEKS assessment (approx 1-2 hours) on Friday, September 9 which will count as Quiz 1. You will need to set up proctoring for this yourself; instructions are in Canvas.
%\end{clist}

\heading{Participation and Recitations}
For most weeks, there will be a check-in activity to ensure that everyone understands the tasks that need to be completed for that week. You will have unlimited attempts for these check-in activities and they will count toward your participation grade. Everyone can (and should) earn 100\% of the points in this category.\\

Math F251X comes with an attached Math F251L Recitation section. There is not a separate recitation Canvas class for the asynchronous class; rather, the recitation activities are built in to the class. For each recitation activity, you will listen to a short video and then you may choose to complete a worksheet that goes with the video. The recitation worksheets are explicitly devoted to bolstering the underlying non-Calculus skills that are nevertheless essential to success in Calculus such as: graphing, algebra, trigonometry, exponential and logarithmic functions, and inverse functions, and they provide targeted instruction on algebra skills that are needed to complete the weekly homework. They also include strategic homework, quiz, and test prep. Ocasionally, there will be a question on one of the check-in activities that is related to the recitation.\\

%Beginning week 3, there will be weekly recitation activities. These are explicitly devoted to bolstering the underlying non-Calculus skills that are nevertheless essential to success in Calculus such as: graphing, algebra, trigonometry, exponential and logarithmic functions, and inverse functions. They also include strategic homework, quiz, and test prep. These can be submitted multiple times so students can (and should) earn 100\% of their recitation points. Note: If you score 80\% or higher on the proctored ALEKS assessment, you may opt out of these recitation worksheets. I recommend completing the worksheets regardless of your ALEKS score.

\heading{Written Homework}
Homework assignments consist of a selection of problems at the end of each section of our textbook. Homework is written (on paper or tablet) and turned in via Gradescope, which may be accessed from Canvas.  Help with scanning homework can be found under \href{https://uaf-math.github.io/calc1/techHelp.html}{Technology Help} on the course webpage. Assignments are usually due on Fridays (by 11:59 PM).  Answers to most problems are provided in the back of the book (or linked from the online text). Complete worked solutions to all problems are provided in advance on Canvas. Thus, your homework will be graded based on \emph{effort} and \emph{completion}. Homework can be turned in up to 24 hours late with no penalty but will not be accepted after that unless there are extenuating circumstances. All students should earn 100\% of their homework points!

The list of homework problems and homework guidelines can be found at the \href{https://uaf-math.github.io/calc1/newhomework.html}{Homework} link on the course webpage. They are also listed in Canvas for each week.

Clearly, it is possible to short-circuit the homework by copying the solutions. It should also be clear that (a) this is a bad idea and (b) your instructor and TA will know you have done this. Our goal in providing answers and solutions is to foster the use of homework as a learning experience. \\

\heading{Weekly Quizzes}
Each week where we have not just had a Midterm or Proficiency, there will be a written quiz in Gradescope. This quiz will test the calculus material that was learned during the previous week. The purpose of the quiz is to help you consolidate your learning from the week, to identify areas of confusion, and then provide an opportunity to fix any misconceptions.  

You will download the quiz from Gradescope and take the quiz on your own without any aids. Once you have attempted all the problems on your own, you will download the solutions to the quiz and \textcolor{red}{correct your quiz in red pen}. You will then upload your corrected quiz. 

\textbf{Your quiz work will earn full points, provided that your initial work is your own and that you fully correct your paper; thus, all students should earn 100\% on the quizzes.} Note that the consequence of this grading scheme is to establish that there is no benefit to your grade by using external sources when initially working the quiz problems and substantial disincentive for violating this rule. \textbf{Specifically, to use external sources when you initially complete the quiz constitutes a violation of academic integrity and the Student Code of Conduct. }

In addition, there is a special proctored ALEKS Quiz that takes place. The ALEKS quiz is is a proctored version of the UAF Math Placement test, which also measures your initial prerequisite knowledge. It counts towards your total Quiz grade and must be completed within the first 3 weeks.

%\heading{Quiz Corrections}
%Beginning with Quiz 2, you can submit revisions for quiz problems you missed. There will be a separate Gradescope assignment, for example ``Quiz 2 Corrections'', which will contain a blank copy of the quiz. On the new copy, write new solutions to any problems that you missed points on, and you will have an opportunity to get points back for correct new solutions. Do not write up problems you got correct originally on the corrected version.

%%to get to compile
\heading{Proctored Assessments}
There are a total of \emph{6 proctored assessments} (one is optional) with dates and testing windows listed below. You will set up the proctoring arrangement through eCampus.
If you live in the Fairbanks area, you can schedule your proctored assessments by going to the eCampus Exam Services site \url{https://ecampus.uaf.edu/student-support/exam-info-students/} and click the yellow box labeled {\tt{Schedule Appointment}} near the middle of the page. If you live outside the Fairbanks area, you should go to eCampus Exam Services site, \url{https://ecampus.uaf.edu/student-support/exam-info-students/} and click the yellow box labeled {\tt{Request Exam}} near the middle of the page.

\begin{center}
\begin{tabular}{| l | l | l |}
\hline
assessment & range of dates & duration \\
\hline \hline
ALEKS Quiz & Monday January 19 -- Friday January 30 & 2 hours\\
\hline
Midterm I & Thursday February 12 or Friday February 13 & 1.5 hours \\
\hline
Derivative Proficiency & Thursday March 5 or Friday March 6 & 1 hour\\
\hline
Derivative Proficiency Retake (Optional) & Thursday March 19 -- Friday March 20 & 1 hour\\
\hline
Midterm II & Thursday April 9 or Friday April 10 & 1.5 hours \\
\hline
%Integral Proficiency& Wed Dec 7 or Thurs Dec 8 & 30 minutes\\
%\hline
Final Exam & Wednesday April 29 -- Friday May 1& 2 hours \\
\hline
\end{tabular}
\end{center}

\heading{Midterms}
There are two midterm exams this semester. Note that the course webpage contains all previous Midterms (with solutions) so a student can know in advance what these are like and has lots of opportunity for practice. The midterms are the same 
for all sections; they are prepared and approved by all instructors teaching the course. 

Make-up midterms will be given only for documented excused absences.\\

\heading{Derivative Proficiency}
The Derivative Proficiency is an assessment covering the routine mechanical skill of differentiation. Note that the course webpage contains all previous derivative proficiencies (with solutions) so a student can know in advance what these are like and has lots of opportunity for practice. The Derivative Proficiency is a proctored assessment and will be graded on a binary scale for each problem (no partial credit). There will be one opportunity to retake the Derivative Proficiency to improve your score.\\

\heading{Final Exam} 
The cumulative final exam will be held at the day/time listed in the
online schedule. Note that the course webpage contains all previous final exams (with solutions) so a student can know in advance what these are like  and has lots of opportunity for practice.
A make-up final exam will be given only in extenuating circumstances, for documented and excused reasons at the discretion of the instructors.

\pagebreak

\heading{Evaluation and Grading Rubric}

%\begin{table}
\begin{multicols}{2}
\begin{tabular}{|c|c|}
\hline
Participation & 3\%\\
\hline
Written Homework & 10\%\\ \hline
Weekly Quizzes& 10\% \\
\hline
Midterm 1 & 20\% \\
\hline
Derivative Proficiency& 12\%\\
\hline
Midterm 2 & 20\%  \\
\hline
Final Exam& 25\% \\
\hline \hline
Total& 100\%\\
\hline
\end{tabular}

\vskip 12pt
Letter grades will be assigned according to the following scale.
This scale is a guarantee; the instructors reserve the right to lower the thresholds if necessary. 

\def\sts{\hskip 0.5em}
\strut\hbox to\hsize{\vbox{\halign{#\hfil\sts&#\hfil\ts&#\hfil\sts&#
\hfil\ts&#\hfil\sts&#\hfil   \cr
A+ & 97--100\% & C+ & 77--79.99\% & F  & $<$ 60\%\cr

A & 93--96.99\% &  C & 70--76.99\%&&\cr
A- & 90--92.99\% & C- & not given&&\cr
B+ & 87--89.99\% & D+ & 67--69.99\%&&\cr
B &  83--86.99\% & D & 63--66.99\%&&\cr
B- & 80-82.99\% & D- & 60--62.99\%&&\cr
}}\hfil}
\end{multicols}

\heading{Office Hours and Communication}
The instructor will host a weekly office hour on Zoom from 3-4pm on Thursdays. Students can also schedule meetings with their instructor outside of regular office hours.\\

\heading{Tutoring and Resources}
\vskip -30pt\strut
\begin{clist}
	\item The Student Success Center on the 6th Floor of the Rasmuson Library, offers drop-in in-person tutoring. 
	See \url{https://www.uaf.edu/dms/mathlab/} for schedules and availability.
	\item One-on-one (or small group) tutoring is available in 
Chapman Building Room 201. You must schedule an
appointment; see \url{https://www.uaf.edu/dms/mathlab/}.
	\item Online tutoring. To make an appointment for online tutoring, go do \url{https://www.uaf.edu/dms/mathlab/}
		\item Student Support Services offers free tutoring in many subjects to students who qualify for their program.
	\item ASUAF offers private tutoring for a small fee (based on student income).
\end{clist}

\vspace{.2in}

\heading{Academic Dishonesty}
Academic dishonesty, including cheating and plagiarism, will not
be tolerated.  It is a violation of the Student Code of Conduct
and will be punished according to UAF procedures.


\pagebreak

\heading{Rules and Policies}
\vskip -20pt


This course is listed as a General Education Math Course. As such this course is expected to meet the 4 general learning outcomes. 

\begin{enumerate}
\item Build knowledge of human institutions, sociocultural processes, and the physical and natural works through the study of mathematics.  Competence will be demonstrated for the foundational information in each subject area, its context and significance, and the methods used in advancing each.

\item Develop intellectual and practical skills across the curriculum, including inquiry and analysis, critical and creative thinking, problem solving, written and oral communication, information literacy, technological competence, and collaborative learning. Proficiency will be demonstrated across the curriculum through critical analysis of proffered information, well-reasoned solutions to problems or inferences drawn from evidence, effective written and oral communication, and satisfactory outcomes of group projects.

\item Acquire tools for effective civic engagement in local through global contexts, including ethical reasoning, intercultural competence, and knowledge of Alaska and Alaska issues.  Facility will be demonstrated through analyses of issues including dimensions of ethics, human and cultural diversity, conflicts and interdependencies, globalization, and sustainability.   

\item Integrate and apply learning, including synthesis and advanced accomplishment across general and specialized studies, adapting them to new settings, questions and responsibilities, and forming a foundation for lifelong learning. Preparation will be demonstrated though production of a a creative or scholarly product that requires broad knowledge, appropriate technical proficiency, information collection, synthesis, interpretation, presentation and reflection.
\end{enumerate}

\subheading{Incomplete Grade} 
Incomplete (I) will only be given in
  DMS courses in cases where
  the student has completed the majority (normally all but the last
  three weeks) of a course with a grade of C or better, but for
  personal reasons beyond his/her control has been unable to complete
  the course during the regular term. Negligence or indifference are
  not acceptable reasons for the granting of an incomplete
  grade. 

\subheading{Late Withdrawals} 
A withdrawal after the deadline
  (currently 9 weeks into the semester) from a DMS course will
  normally be granted only in cases where the student is performing
  satisfactorily (i.e., C or better) in a course, but has exceptional
  reasons, beyond his/her control, for being unable to complete the
  course. These exceptional reasons should be detailed in writing to
  the instructor, department head and dean.

\subheading{No Early Final Examinations}
Final examinations for DMS
  courses shall not be held earlier than the date and time published
  in the official term schedule. Normally, a student will not be
  allowed to take a final exam early. Exceptions can be made by
  individual instructors, but should only be allowed in exceptional
  circumstances and in a manner which doesn't endanger the security of
  the exam.

\begin{center}
\textbf{\large{Official UAF Syllabus Addendum}}
\end{center}


\noindent{\bf Student protections statement:} The university respects and upholds the principles of due process and a fair and equitable process as specified in the Board of Regents' Policy 09.02 Student Rights and Responsibilities. For more information regarding the rights and responsibilities of students, refer to the Office of Rights, Compliance and Accountability website. You are encouraged to read the Board of Regents' policy carefully to fully understand your responsibilities to our community.

We strive to create a safe and respectful environment for all members of our community. If you have questions about expectations of you as a student or believe your rights are being violated, we encourage you to reach out to the  Office of Rights, Compliance and Accountability for help. UAF reserves the right to suspend, expel or take other necessary and appropriate action in cases where a student is unable or unwilling to uphold community standards and campus safety.

For more information on your rights as a student and the resources available to you to resolve problems, please go to the following site:\\ \url{https://catalog.uaf.edu/academics-regulations/students-rights-responsibilities/}

\noindent{\bf Disability services statement:} I will work with the Office of Disability Services to provide reasonable accommodation to students with disabilities.

\noindent{\bf ASUAF advocacy statement:} The Associated Students of the University of Alaska Fairbanks, the student government of UAF, offers advocacy services to students who feel they are facing issues with staff, faculty, and/or other students specifically if these issues are hindering the ability of the student to succeed in their academics or go about their lives at the university. Students who wish to utilize these services can contact the Student Advocacy Director by visiting the ASUAF office or emailing asuaf.office@alaska.edu. 



\noindent{\bf Student Academic Support:}
\begin{itemize}
\setlength\itemsep{0em}
        \item Communication Center (907-474-7007, \mailto{uaf-commcenter@alaska.edu}, Student Success Center, 6th Floor Room 677 Rasmuson Library)
        \item Writing Center (907-474-5314, \mailto{uaf-writing-center@alaska.edu}, Student Success Center, 6th Floor Room 677 Rasmuson Library)
\item UAF Math Services (907-474-7332, \mailto{uaf-traccloud@alaska.edu})


\begin{itemize}
\item Drop-in tutoring, Student Success Center, 6th Floor Room 672 Rasmuson Library

\item 1:1 tutoring (by appointment only), 6th Floor Room 677 Rasmuson Library

\item Online tutoring (by appointment only) available

https://www.uaf.edu/dms/mathlab/, available at the Student Success Center
\end{itemize}

\item Developmental Math Lab, Gruening 406
\item The Debbie Moses Learning Center at CTC (907-455-2860, 604 Barnette St, Room 120,\\ \url{https://www.ctc.uaf.edu/student-services/student-success-center/})
\item For more information and resources, please see the Academic Advising Resource List (\url{https://www.uaf.edu/advising/students/index.php})
\end{itemize}

\noindent{\bf Student Resources:}
\begin{itemize}
\setlength\itemsep{0em}
\item Disability Services (907-474-5655, \mailto{uaf-disability-services@alaska.edu}, 110 Eielson Building)
\item Student Health \& Counseling [free counseling sessions available] (907-474-7043, \url{https://www.uaf.edu/chc/appointments.php}, Whitaker Building, Room 206, Health, Safety \& Security Bldg --- same building as Fire and Police)
\item Office of Rights, Compliance and Accountability (907-474-7300, \mailto{uaf-orca@alaska.edu}, 3rd Floor, Constitution Hall)
\item Associated Students of the University of Alaska Fairbanks (ASUAF) or ASUAF Student Government (907-474-7355, \mailto{asuaf.office@alaska.edu}{asuaf.office@alaska.edu}, Wood Center 119)
\end{itemize}

\noindent{\bf Nondiscrimination statement:}
Nondiscrimination statement: The University of Alaska is an equal opportunity/equal access employer, educational institution and provider. The University of Alaska does not discriminate on the basis of race, religion, color, national origin, citizenship, age, sex, physical or mental disability, status as a protected veteran, marital status, changes in marital status, pregnancy, childbirth or related medical conditions, parenthood, sexual orientation, gender identity, political affiliation or belief, genetic information, or other legally protected status. The University's commitment to nondiscrimination, including against sex discrimination, applies to students, employees, and applicants for admission and employment. Contact information, applicable laws, and complaint procedures are included on UA's statement of nondiscrimination available at \url{www.alaska.edu/nondiscrimination}.

\begin{tabular}{l}
UAF Office of Rights, Compliance and Accountability\\
1692 Tok Lane\\
3rd floor, Constitution Hall, Fairbanks, AK 99775\\
907-474-7300\\
\url{uaf-orca@alaska.edu}
\end{tabular}


 \scriptsize syllabus version: \today \normalsize

\end{document}

